% !TeX root = ../thuthesis-example.tex

\begin{resume}

  \section*{个人简历}

  1984 年 12 月 8 日出生于山东省东营市。

  2003年 年 9 月考入哈尔滨工业大学自动化专业，2007 年 7 月本科毕业并获得工学学士学位。

  2007年 9 月免试进入哈尔滨工业大学控制理论与控制工程专业，2009 年 7 月硕士毕业并获得工学硕士学位。
  
  2009年 9 月考入中国科学院自动化研究所模式识别专业，2012 年 7 月博士毕业并获得工学博士学位。

  2023 年 9 月考入清华大学经管学院工商管理专业攻读专业硕士至今。


  \section*{在学期间完成的相关学术成果}

  \subsection*{学术论文}

 % \begin{achievements}
%  \item Yang Y, Ren T L, Zhang L T, et al. Miniature microphone with silicon-based ferroelectric thin films[J]. Integrated Ferroelectrics, 2003, 52:229-235.
%  \item 杨轶, 张宁欣, 任天令, 等. 硅基铁电微声学器件中薄膜残余应力的研究[J]. 中国机械工程, 2005, 16(14):1289-1291.
%  \item 杨轶, 张宁欣, 任天令, 等. 集成铁电器件中的关键工艺研究[J]. 仪器仪表学报, 2003, 24(S4):192-193.
%  \item Yang Y, Ren T L, Zhu Y P, et al. PMUTs for handwriting recognition. In press[J]. (已被Integrated Ferroelectrics录用)
%  \end{achievements}


  \subsection*{专利}

%  \begin{achievements}
%    \item 任天令, 杨轶, 朱一平, 等. 硅基铁电微声学传感器畴极化区域控制和电极连接的方法: 中国, CN1602118A[P]. 2005-03-30.
%    \item Ren T L, Yang Y, Zhu Y P, et al. Piezoelectric micro acoustic sensor based on ferroelectric materials: USA, No.11/215, 102[P]. (美国发明专利申请号.)
%  \end{achievements}

\end{resume}



% 本科生格式：

% \begin{resume}
%   \section*{学术论文}
%
%   \begin{achievements}
%     \item ZHOU R, HU C, OU T, et al. Intelligent GRU-RIC Position-Loop
%       Feedforward Compensation Control Method with Application to an
%       Ultraprecision Motion Stage[J], IEEE Transactions on Industrial
%       Informatics, 2024, 20(4): 5609-5621.
%
%     \item 杨轶, 张宁欣, 任天令, 等. 硅基铁电微声学器件中薄膜残余应力的研究[J].
%       中国机械工程, 2005, 16(14):1289-1291.
%
%     \item YANG Y, REN T L, ZHU Y P, et al. PMUTs for handwriting recognition.
%       In press[J]. (已被Integrated Ferroelectrics录用)
%
%   \end{achievements}
%
%
%   \section*{专利}
%
%   \begin{achievements}
%     \item 胡楚雄, 付宏, 朱煜, 等. 一种磁悬浮平面电机: ZL202011322520.6[P]. 2022-04-01.
%
%     \item REN T L, YANG Y, ZHU Y P, et al. Piezoelectric micro acoustic sensor
%       based on ferroelectric materials: No.11/215, 102[P]. (美国发明专利申请号.)
%
%   \end{achievements}
% \end{resume}
